\par\indent Bekanntlich haben unabhängig von \textsc{Gauss} und voneinander auch \textsc{J. R.
Argand} (1806) und \textsc{C. Wessel} (1799) die geometrische Versinnlichung der
komplexen Grössen gefunden; neuerdings gründet man die Berechtigung, komplexe
Grössen als Substrate für die formalen Rechenoperationen zu betrachten,
auf die von \textsc{W. R. Hamilton}\footnote{\textsc{R. W. Hamilton}, Transactions of the Royal Irish Academy 17 (Dublin 1837), S. 293 und \textit{Lectures on Quaternions} (ibid. 1853) Vorrede.} und \textsc{Johann Bolyai}\footnote{\textsc{Johann Bolyai}, \textit{Responsio ad quaestionem .\ .\ .\ . ab Incl. Soc. Jablonowskiana Lipsiae anno
1837 motam}, zum ersten Male herausgegeben von \textsc{P. Stäckel}, Mathem.-Naturw. Berichte aus Ungarn 16
(1899), S. 281, deutsch in \textsc{Stäckel}, \textit{Die beiden Bolyai} II, 1913, S. 221.} herrührende Auffassung\par
\clearpage\SourceLeaf{141}\ArithmeticHead{56}{L}{L. SCHLESINGER, ÜBER GAUSS’ FUNKTIONENTHEORETISCHE ARBEITEN.}
\noindent dieser Grössen als Paare von reellen bezw. Quadrupel von positiven Grössen.
\textsc{Johann Bolyai} hat in seiner Abhandlung auch eine Kritik der von \textsc{Gauss}
entwickelten Auffassung gegeben, aus der wir nur eine Stelle wiedergeben
wollen, da die andern Einwände sich als gegenstandslos erweisen\footnote{Der zweite Einwand \textsc{Bolyais} (\textsc{Stäckel} a. a. O., S. 233): »Man begreift nicht, wie man zu dem
Schlusse gelangt, dass \(+i\) (wie auch \(-i\)) die mittlere Proportionale zwischen \(+1\) und \(-1\) ist, umsomehr
als vorher die Proportion nicht allgemein erklärt worden ist und auch Rhomben statt der Quadrate genommen
werden können« zeigt, dass \textsc{Johann Bolyai} bei seiner Kritik nur die \textit{Anzeige}, nicht die Abhandlung
selbst vor Augen hatte, da ja, wie wir gesehen haben, in der letzteren die allgemeine Definition der
Proportion tatsächlich gegeben ist, während in der \textit{Anzeige} nur der Übergang von \(\pm1\) zu \(\pm i\) erörtert
wird. Dass die quadratische Anordnung nur eine \textls[120]{Möglichkeit} verkörpert, von der abzugehn »kein Grund
vorhanden ist« wird in der \textit{Anzeige} ausdrücklich betont. — Trotzdem hat \textsc{Stäckel} (a. a. O. I, S. 133) und
ebenso \textsc{Fraenkel}, \textit{Zahlbegriff und Algebra bei Gauss}, \textit{Materialien} VIII, 1920, S. 42, diesen Einwand
\textsc{Bolyais} für berechtigt erklärt.}. \textsc{Bolyai}
sagt in bezug auf die Darstellung der \textit{Anzeige}: »Die vorliegende Auseinandersetzung
stützt sich .\ .\ .\ . auf die Betrachtung des Raumes, die man in der
Arithmetik vermeiden soll« (\textsc{Stäckel}, a. a. O., S. 233). Natürlich hat \textsc{Gauss}
von diesem Einwand \textsc{Bolyais} niemals Kenntnis erhalten, umso bemerkenswerter
ist es, dass er in verschiedenen Äusserungen aus späterer Zeit ähnliche Bedenken
angedeutet und eine rein abstrakte, von allem Räumlichen unabhängige
Rechtfertigung der komplexen Grössen erstrebt hat. So sagt er in dem erwähnten
Brief an \textsc{Hansen} (Werke XII, S. 8): »Der wahre Sinn des \(\sqrt{-1}\)
steht mir dabei mit grosser Lebendigkeit vor der Seele, aber es wird sehr
schwer sein, ihn in Worte zu fassen, die immer nur ein vages, in der Luft
schwebendes Bild geben können .\ .\ .«, ferner in dem Brief an \textsc{Drobisch}\footnote{Der Mathematiker und Philosoph \textsc{M. W. Drobisch} (1802—1896) war übrigens derjenige, der die
von \textsc{Johann Bolyai} bearbeitete Preisaufgabe gestellt hatte, in deren Thema auf \textsc{Gauss’} Versinnlichung
der imaginären Grössen ausdrücklich hingewiesen wird. Näheres bei \textsc{Stäckel}, a. a. O. I, S. 125 ff.}
(Werke X, 1, S. 106): »Nur ist die Darstellung der imaginären Grössen in den
Relationen der Punkte in plano nicht sowohl ihr Wesen selbst, welches höher
und allgemeiner aufgefasst werden muss, als vielmehr das uns Menschen reinste,
oder vielleicht einzig ganz reine Beispiel ihrer Anwendung«, endlich heisst es
in der Jubiläumsschrift (1849, Werke III, S. 79) mit Beziehung auf den Beweis
der Wurzelexistenz: »Im Grunde gehört aber der eigentliche Inhalt der
ganzen Argumentation einem höheren, von Räumlichem unabhängigen Gebiet\par
\clearpage\SourceLeaf{142}\ArithmeticHead{57}{R}{III e. DIE INAUGURALDISSERTATION.}
\noindent der allgemeinen abstrakten Grössenlehre an, dessen Gegenstand die nach der
Stetigkeit zusammenhängenden Grössenkombinationen sind, einem Gebiete, welches
zur Zeit noch wenig angebaut ist, und in welchem man sich auch nicht
bewegen kann ohne eine von räumlichen Bildern entlehnte Sprache«, vergl.
auch Werke X, 1, S. 408, wo auf diese Stelle der Jubiläumsschrift hingedeutet
wird. Namentlich aus der letztzitierten Stelle geht hervor, dass \textsc{Gauss} in den
»Grössenkombinationen«, d. h. wohl in der Betrachtung von Systemen reeller
Grössen, das Substrat der allgemeinen abstrakten Grössenlehre gesehen hat;
in dieser Richtung liegen die aus dem Nachlass veröffentlichten Untersuchungen
über drei- und viergliedrige komplexe Grössen (Werke VIII, S. 353—362,
vergl. \textsc{Stäckel}, Abh. 4 dieses Bandes, Nr. 20), ferner haben wir in dem
Nachlassstück (Werke X, 1, S. 408) in den Bezeichnungen »Stelle, Zug, Schicht«
für eine einzelne komplexe Grösse, bezw. für ein- und zweidimensionale »nach
der Stetigkeit zusammenhängende« Reihen solcher Grössen ein Spezimen dafür,
wie sich \textsc{Gauss} die im »Felde der abstrakten komplexen Grössen« zu verwendende
Sprache gedacht hat.

Die schon in der \textit{Anzeige} ausgesprochene Absicht, »den Gegenstand vollständiger
zu bearbeiten, wo dann auch die Frage, warum die Relationen zwischen
Dingen, die eine Mannigfaltigkeit von mehr als zwei Dimensionen darbieten,
nicht noch andere, in der allgemeinen Arithmetik zulässige Arten von
Grössen liefern können, ihre Beantwortung finden wird« hat \textsc{Gauss} leider nicht
ausgeführt.
